%═══════════════════════════════════════════════════════════════════════════════
% APÉNDICE A — Herramientas de Análisis Numérico Computacional
%
% Objetivo: proporcionar las definiciones, métricas y herramientas algebraicas
% que se usan de forma transversal en todas las prácticas. Este apéndice es
% autocontenido y actúa como referencia rápida; cada práctica lo cita cuando
% introduce métricas de error, número de condición o convergencia empírica.
%═══════════════════════════════════════════════════════════════════════════════
\chapter{Herramientas de Análisis Numérico Computacional}
\label{ap:herramientas}

\begin{flushright}
\itshape
``A weapon is never the right tool for the job.\\
Neither is a number without an error bar.''\\[2pt]
\upshape --- Richard W. Hamming (parafraseado)
\end{flushright}
\bigskip

Este apéndice consolida las herramientas matemáticas que aparecen de forma
recurrente en las prácticas: normas vectoriales y matriciales, métricas de
error computacional, número de condición, descomposición en valores singulares
(SVD) y estimación empírica de la tasa de convergencia. No es un texto de
álgebra lineal numérica; para ese fin véanse \cite{golub2013matrix} y
\cite{higham2002accuracy}. El objetivo es que el estudiante tenga a mano
las fórmulas exactas y su interpretación física, sin necesidad de buscarlas
en otra fuente durante el trabajo experimental.

%═══════════════════════════════════════════════════════════════════════════════
\section{Normas vectoriales y matriciales}
\label{ap:normas}
%═══════════════════════════════════════════════════════════════════════════════

\subsection{Normas en \texorpdfstring{$\mathbb{R}^n$}{Rn}}
\label{ap:normas_vec}

\begin{definicion}{Norma vectorial $\ell^p$}
\label{def:norma_lp}
Sea $\mathbf{v} \in \mathbb{R}^n$. Para $1 \leq p < \infty$, la
\textbf{norma $\ell^p$} se define como
\[
  \|\mathbf{v}\|_p \;=\; \left(\sum_{i=1}^{n} |v_i|^p\right)^{1/p},
\]
y la \textbf{norma $\ell^\infty$} (norma supremo o norma máximo) como
\[
  \|\mathbf{v}\|_\infty \;=\; \max_{1 \leq i \leq n} |v_i|.
\]
Los casos de uso más frecuente en física computacional son $p=1$, $p=2$
(norma euclídea) y $p=\infty$.
\end{definicion}

Las tres normas más comunes tienen interpretaciones físicas directas: la
norma $\ell^1$ mide la distancia de Manhattan (suma de desplazamientos
axiales), la norma $\ell^2$ es la distancia euclídea usual, y la norma
$\ell^\infty$ captura el peor caso puntual —la desviación máxima entre
dos vectores.

\begin{proposicion}{Equivalencia de normas en dimensión finita}
\label{prop:equiv_normas}
En $\mathbb{R}^n$ con $n < \infty$, todas las normas son equivalentes: para
cualesquiera normas $\|\cdot\|_\alpha$ y $\|\cdot\|_\beta$ existen
constantes $0 < c \leq C < \infty$ tales que
\[
  c\,\|\mathbf{v}\|_\alpha \;\leq\; \|\mathbf{v}\|_\beta \;\leq\;
  C\,\|\mathbf{v}\|_\alpha \quad \forall\, \mathbf{v} \in \mathbb{R}^n.
\]
En particular, para las normas $\ell^1$, $\ell^2$ y $\ell^\infty$:
\[
  \|\mathbf{v}\|_\infty \;\leq\; \|\mathbf{v}\|_2 \;\leq\;
  \|\mathbf{v}\|_1 \;\leq\; \sqrt{n}\,\|\mathbf{v}\|_2 \;\leq\;
  n\,\|\mathbf{v}\|_\infty.
\]
\end{proposicion}

\begin{observacion}{Relevancia en la práctica}
La equivalencia garantiza que la elección de norma no altera
cualitativamente la noción de convergencia. Sin embargo, los factores
$\sqrt{n}$ y $n$ son numéricamente significativos cuando $n$ es grande:
una norma $\ell^1$ pequeña no implica que la norma $\ell^\infty$ sea
pequeña a menos que el error sea uniforme sobre todas las componentes.
\end{observacion}

\subsection{Normas en espacios funcionales}
\label{ap:normas_func}

\begin{definicion}{Normas $L^p$ continuas}
\label{def:norma_Lp}
Sea $f : [a,b] \to \mathbb{R}$ una función medible. Las normas $L^p$
continuas se definen como
\begin{align}
  \|f\|_{L^1} &= \int_a^b |f(x)|\,dx, \label{eq:norm_L1_cont}\\[4pt]
  \|f\|_{L^2} &= \left(\int_a^b |f(x)|^2\,dx\right)^{1/2}, \label{eq:norm_L2_cont}\\[4pt]
  \|f\|_{L^\infty} &= \operatorname{ess\,sup}_{x \in [a,b]} |f(x)|. \label{eq:norm_Linf_cont}
\end{align}
El espacio $L^2[a,b]$ dotado del producto interno
$\langle f,g\rangle = \int_a^b f(x)\,g(x)\,dx$
es un espacio de Hilbert \cite{reed1980functional}.
\end{definicion}

\begin{fisicaconexion}{Funciones como vectores de dimensión infinita}
El espacio $L^2[a,b]$ es el análogo de dimensión infinita de $\mathbb{R}^n$.
Los elementos de una base ortonormal $\{\phi_j\}_{j=0}^\infty$ satisfacen
$\langle \phi_j, \phi_k \rangle = \delta_{jk}$, exactamente como los
vectores de la base canónica en $\mathbb{R}^n$. Truncar una serie de Fourier
a $p$ términos es proyectar el vector $f \in L^2$ sobre el subespacio
generado por los primeros $p$ vectores base; el teorema de la mejor
aproximación garantiza que esta proyección minimiza la distancia $L^2$.
En las prácticas, las versiones discretas de estas normas (evaluadas sobre
$n$ puntos) serán las métricas de error operativas.
\end{fisicaconexion}

\subsection{Normas matriciales inducidas}
\label{ap:normas_mat}

\begin{definicion}{Norma matricial inducida}
\label{def:norma_mat_ind}
Dada una norma vectorial $\|\cdot\|_p$ en $\mathbb{R}^n$, la norma matricial
inducida (u \textit{operador}) de $\mathbf{A} \in \mathbb{R}^{m \times n}$ es
\[
  \|\mathbf{A}\|_p \;=\; \sup_{\mathbf{v} \neq \mathbf{0}}
  \frac{\|\mathbf{A}\mathbf{v}\|_p}{\|\mathbf{v}\|_p}
  \;=\; \max_{\|\mathbf{v}\|_p = 1} \|\mathbf{A}\mathbf{v}\|_p.
\]
\end{definicion}

Los casos más utilizados son:
\[
  \|\mathbf{A}\|_1 = \max_{1 \leq j \leq n}\sum_{i=1}^m |A_{ij}|
  \quad \text{(máximo de sumas de columnas)},
\]
\[
  \|\mathbf{A}\|_\infty = \max_{1 \leq i \leq m}\sum_{j=1}^n |A_{ij}|
  \quad \text{(máximo de sumas de filas)},
\]
\[
  \|\mathbf{A}\|_2 = \sigma_{\max}(\mathbf{A})
  \quad \text{(mayor valor singular, norma espectral)}.
\]

La norma de Frobenius, aunque no es norma inducida, aparece frecuentemente:
\[
  \|\mathbf{A}\|_F = \left(\sum_{i,j} A_{ij}^2\right)^{1/2}
  = \left(\sum_k \sigma_k^2\right)^{1/2},
\]
donde $\sigma_k$ son los valores singulares de $\mathbf{A}$.

%═══════════════════════════════════════════════════════════════════════════════
\section{Métricas de error en experimentos computacionales}
\label{ap:metricas_error}
%═══════════════════════════════════════════════════════════════════════════════

En física experimental, ninguna medición se reporta sin su barra de error.
El mismo principio aplica en física computacional: cualquier resultado
numérico debe acompañarse de una métrica de discrepancia respecto a la
solución de referencia. A continuación se definen las métricas estándar
utilizadas en este laboratorio.

\begin{definicion}{Error absoluto y error relativo (puntuales)}
\label{def:error_puntual}
Para un valor calculado $\hat{v}$ y su referencia $v_{\mathrm{ref}}$, el
\textbf{error absoluto} y el \textbf{error relativo} puntuales son:
\[
  e_{\mathrm{abs}} = |\hat{v} - v_{\mathrm{ref}}|,
  \qquad
  e_{\mathrm{rel}} = \frac{|\hat{v} - v_{\mathrm{ref}}|}{|v_{\mathrm{ref}}|}
  \quad (v_{\mathrm{ref}} \neq 0).
\]
\end{definicion}

\begin{definicion}{Métricas de error global discretas}
\label{def:metricas_discretas}
Dados $n$ pares de valores calculados y de referencia
$\{(\hat{y}_i, y_i)\}_{i=1}^n$, se definen:

\begin{itemize}[nosep, leftmargin=1.5em]
  \item \textbf{Error medio absoluto} (\textit{mean absolute error}, MAE):
  \[
    \mathrm{MAE} \;=\; \frac{1}{n}\sum_{i=1}^n |\hat{y}_i - y_i|
    \;=\; \frac{1}{n}\|\hat{\mathbf{y}} - \mathbf{y}\|_1.
  \]
  \item \textbf{Error cuadrático medio} (\textit{root mean square error}, RMSE):
  \[
    \mathrm{RMSE} \;=\; \sqrt{\frac{1}{n}\sum_{i=1}^n (\hat{y}_i - y_i)^2}
    \;=\; \frac{1}{\sqrt{n}}\|\hat{\mathbf{y}} - \mathbf{y}\|_2.
  \]
  \item \textbf{Error máximo} (norma $\ell^\infty$ del residuo):
  \[
    E_\infty \;=\; \max_{1 \leq i \leq n}|\hat{y}_i - y_i|
    \;=\; \|\hat{\mathbf{y}} - \mathbf{y}\|_\infty.
  \]
  \item \textbf{Norma $\ell^2$ del residuo} (sin normalizar):
  \[
    E_2 \;=\; \|\hat{\mathbf{y}} - \mathbf{y}\|_2
    \;=\; \sqrt{\sum_{i=1}^n (\hat{y}_i - y_i)^2}.
  \]
\end{itemize}
\end{definicion}

\begin{observacion}{Escala y comparación de métricas}
El RMSE y el MAE tienen las mismas unidades que $y$ y son directamente
comparables entre experimentos con el mismo $n$. La norma $E_\infty$ es
la más sensible a valores atípicos (outliers) porque captura el peor caso
puntual. En aproximación de funciones, $E_\infty$ es la métrica relevante
para garantías uniformes (tipo Chebyshev); el RMSE es más adecuado cuando
los errores son gaussianos. Las normas $E_2$ y RMSE solo difieren por el
factor $1/\sqrt{n}$; reportar RMSE en lugar de $E_2$ permite comparar
experimentos con distinto número de puntos.
\end{observacion}

La siguiente tabla resume todas las métricas con sus fórmulas compactas en
notación vectorial:

\begin{table}[htbp]
\centering
\caption{Métricas de error computacional y sus equivalentes en normas vectoriales.
  $\mathbf{r} = \hat{\mathbf{y}} - \mathbf{y}$ es el vector de residuos.}
\label{tab:metricas}
\begin{tabular}{llll}
  \toprule
  \textbf{Métrica} & \textbf{Símbolo} & \textbf{Fórmula explícita}
    & \textbf{Norma vectorial} \\
  \midrule
  Error medio absoluto      & MAE        & $\tfrac{1}{n}\sum|r_i|$
    & $\tfrac{1}{n}\|\mathbf{r}\|_1$ \\[4pt]
  Error cuadrático medio    & RMSE       & $\sqrt{\tfrac{1}{n}\sum r_i^2}$
    & $\tfrac{1}{\sqrt{n}}\|\mathbf{r}\|_2$ \\[4pt]
  Norma $\ell^2$ residuo    & $E_2$      & $\sqrt{\sum r_i^2}$
    & $\|\mathbf{r}\|_2$ \\[4pt]
  Error máximo              & $E_\infty$ & $\max_i|r_i|$
    & $\|\mathbf{r}\|_\infty$ \\[4pt]
  Error relativo $\ell^2$   & $E_{\mathrm{rel}}$ &
    $\|\mathbf{r}\|_2 / \|\mathbf{y}\|_2$ & --- \\
  \bottomrule
\end{tabular}
\end{table}

\FloatBarrier

%═══════════════════════════════════════════════════════════════════════════════
\section{Número de condición y estabilidad de sistemas lineales}
\label{ap:condicion}
%═══════════════════════════════════════════════════════════════════════════════

El número de condición es quizá el observable numérico más importante de
un sistema lineal: mide cuánto se amplifican los errores de entrada (ruido
en los datos, errores de redondeo) en la solución de salida.

\begin{definicion}{Número de condición de una matriz}
\label{def:numero_condicion}
Para una matriz cuadrada e invertible $\mathbf{A} \in \mathbb{R}^{n \times n}$,
el \textbf{número de condición} en la norma $p$ es
\[
  \kappa_p(\mathbf{A}) \;=\; \|\mathbf{A}\|_p \cdot \|\mathbf{A}^{-1}\|_p.
\]
En la norma espectral ($p=2$), se obtiene el número de condición más común:
\[
  \kappa_2(\mathbf{A}) \;=\; \frac{\sigma_{\max}}{\sigma_{\min}},
\]
donde $\sigma_{\max}$ y $\sigma_{\min}$ son el mayor y menor valor singular de
$\mathbf{A}$, respectivamente. Para una matriz rectangular
$\mathbf{A} \in \mathbb{R}^{m \times n}$ con $m \geq n$, se usa la misma
expresión con los valores singulares no nulos.
\end{definicion}

\begin{teorema}{Amplificación de perturbaciones en sistemas lineales}
\label{thm:perturbacion}
Considérese el sistema $\mathbf{A}\mathbf{x} = \mathbf{b}$ y su versión
perturbada $\mathbf{A}(\mathbf{x} + \delta\mathbf{x}) = \mathbf{b} +
\delta\mathbf{b}$. Entonces
\[
  \frac{\|\delta\mathbf{x}\|}{\|\mathbf{x}\|}
  \;\leq\; \kappa(\mathbf{A})\,\frac{\|\delta\mathbf{b}\|}{\|\mathbf{b}\|}.
\]
El número de condición $\kappa(\mathbf{A})$ es la cota superior para la
razón de amplificación del error relativo. Analogamente, si $\mathbf{A}$ se
perturba en $\delta\mathbf{A}$:
\[
  \frac{\|\delta\mathbf{x}\|}{\|\mathbf{x}\|}
  \;\leq\; \kappa(\mathbf{A})\,\frac{\|\delta\mathbf{A}\|}{\|\mathbf{A}\|}.
\]
\end{teorema}

\begin{observacion}{Interpretación logarítmica del número de condición}
Si $\kappa(\mathbf{A}) \approx 10^k$ y se trabaja con aritmética de punto
flotante de precisión $\varepsilon_{\mathrm{mach}} \approx 10^{-16}$
(doble precisión), el número de dígitos decimales de precisión que se
\emph{pierden} en la solución es aproximadamente $k$. Dicho de otro modo:
se recuperan del orden de $16 - k$ dígitos significativos. Un sistema con
$\kappa \approx 10^{12}$ proporciona soluciones con tan solo $\sim 4$
dígitos de precisión en doble precisión.
\end{observacion}

\begin{fisicaconexion}{Número de condición como razón de escalas}
En mecánica estructural, $\kappa(\mathbf{K})$ para la matriz de rigidez
de un sólido mide la razón entre la rigidez del modo más rígido y el más
blando: un sistema compuesto por materiales muy dispares (e.g., acero y
goma) tiene $\kappa \gg 1$. En redes de circuitos, $\kappa(\mathbf{G})$
para la matriz de conductancias mide la razón entre la resistencia mínima
y máxima. En ambos casos, un $\kappa$ grande indica la coexistencia de
escalas muy dispares, que es exactamente lo que hace numéricamente
inestable la resolución del sistema. El oscilador armónico amortiguado con
frecuencia muy alta y amortiguamiento muy bajo produce una matriz de
evolución bien condicionada; agregar términos polinomiales de grado alto a
la base puede llevar el número de condición a $\sim 10^{20}$, haciendo la
inversión numéricamente inútil en doble precisión.
\end{fisicaconexion}

\subsection{Número de condición de sistemas sobredeterminados}
\label{ap:cond_sobred}

En problemas de mínimos cuadrados, el sistema a resolver es
$\min_{\boldsymbol{\theta}} \|\mathbf{A}\boldsymbol{\theta} - \mathbf{y}\|_2^2$
con $\mathbf{A} \in \mathbb{R}^{n \times p}$, $n > p$.
Las ecuaciones normales $\mathbf{A}^\top\mathbf{A}\,\boldsymbol{\theta} =
\mathbf{A}^\top\mathbf{y}$ tienen número de condición
\[
  \kappa_2(\mathbf{A}^\top\mathbf{A})
  \;=\; \kappa_2(\mathbf{A})^2.
\]
Este cuadrado es la razón por la que se prefiere resolver el sistema
sobredeterminado directamente mediante SVD o factorización QR, en lugar de
construir explícitamente las ecuaciones normales: la formación de
$\mathbf{A}^\top\mathbf{A}$ duplica el número de dígitos significativos que
se pierden.

%═══════════════════════════════════════════════════════════════════════════════
\section{Aritmética de punto flotante y épsilon de máquina}
\label{ap:floating_point}
%═══════════════════════════════════════════════════════════════════════════════

\begin{definicion}{Épsilon de máquina}
\label{def:epsilon_mach}
El \textbf{épsilon de máquina} $\varepsilon_{\mathrm{mach}}$ es el menor
número positivo de punto flotante tal que $1 + \varepsilon_{\mathrm{mach}}
> 1$ en aritmética de precisión finita. En el estándar IEEE\,754:
\[
  \varepsilon_{\mathrm{mach}} \;=\;
  \begin{cases}
    2^{-23} \approx 1.19 \times 10^{-7} & \text{precisión simple (32 bits),} \\[3pt]
    2^{-52} \approx 2.22 \times 10^{-16} & \text{precisión doble (64 bits).}
  \end{cases}
\]
\end{definicion}

El épsilon de máquina define el piso de error numérico. Toda operación
introduce un error relativo del orden de $\varepsilon_{\mathrm{mach}}$;
algoritmos que suman $n$ términos acumulan errores del orden de
$n\varepsilon_{\mathrm{mach}}$ (cancelación benigna) o, en el peor caso,
sufren \emph{cancelación catastrófica}: la resta de dos números casi iguales
puede resultar en un valor con casi ningún dígito correcto.

\begin{ejemplo}{Cancelación catastrófica en la fórmula cuadrática}
La fórmula $x_{1,2} = (-b \pm \sqrt{b^2 - 4ac}) / (2a)$ sufre cancelación
catastrófica cuando $b^2 \gg 4ac$ (raíces muy dispares): la resta
$-b + \sqrt{b^2-4ac}$ cancela casi todos los dígitos significativos. La
reformulación estable es calcular la raíz de mayor módulo con la fórmula
directa y obtener la otra como $x_2 = c / (a\, x_1)$ usando la relación de
Vieta.
\end{ejemplo}

%═══════════════════════════════════════════════════════════════════════════════
\section{Descomposición en valores singulares (SVD)}
\label{ap:svd}
%═══════════════════════════════════════════════════════════════════════════════

La SVD es la herramienta central del análisis de datos lineales y de la
solución estable de sistemas sobredeterminados. Aparece explícitamente en
la segunda práctica, pero ya en la primera es necesaria para calcular
$\kappa(\mathbf{A})$ y como alternativa estable a las ecuaciones normales.

\begin{teorema}{Descomposición en valores singulares}
\label{thm:svd}
Toda matriz $\mathbf{A} \in \mathbb{R}^{m \times n}$ con $m \geq n$ admite
una factorización
\[
  \mathbf{A} \;=\; \mathbf{U}\,\boldsymbol{\Sigma}\,\mathbf{V}^\top,
\]
donde $\mathbf{U} \in \mathbb{R}^{m \times m}$ y
$\mathbf{V} \in \mathbb{R}^{n \times n}$ son matrices ortogonales
($\mathbf{U}^\top\mathbf{U} = \mathbf{I}_m$,
$\mathbf{V}^\top\mathbf{V} = \mathbf{I}_n$), y
$\boldsymbol{\Sigma} \in \mathbb{R}^{m \times n}$ es diagonal con entradas
no negativas $\sigma_1 \geq \sigma_2 \geq \cdots \geq \sigma_n \geq 0$
llamadas \textbf{valores singulares} de $\mathbf{A}$. La factorización es
única (salvo signo de las columnas de $\mathbf{U}$ y $\mathbf{V}$) cuando
los valores singulares son distintos.
\end{teorema}

Las columnas de $\mathbf{U}$ (vectores singulares izquierdos) forman una
base ortonormal para $\mathbb{R}^m$ y las columnas de $\mathbf{V}$ (vectores
singulares derechos) para $\mathbb{R}^n$. La SVD ``económica'' o ``thin
SVD'' retiene solo las primeras $r = \mathrm{rank}(\mathbf{A})$ columnas
de $\mathbf{U}$, dando $\mathbf{A} = \hat{\mathbf{U}}\hat{\boldsymbol{\Sigma}}
\mathbf{V}^\top$ con $\hat{\mathbf{U}} \in \mathbb{R}^{m\times r}$ y
$\hat{\boldsymbol{\Sigma}} \in \mathbb{R}^{r \times r}$.

\begin{proposicion}{Pseudoinversa de Moore-Penrose vía SVD}
\label{prop:pseudoinversa}
La solución de mínimos cuadrados de mínima norma del sistema
$\mathbf{A}\boldsymbol{\theta} = \mathbf{y}$ es
$\boldsymbol{\theta}^+ = \mathbf{A}^+\mathbf{y}$, donde la
\textbf{pseudoinversa} $\mathbf{A}^+$ se obtiene como
\[
  \mathbf{A}^+ \;=\; \mathbf{V}\,\boldsymbol{\Sigma}^+\,\mathbf{U}^\top,
\]
con $\boldsymbol{\Sigma}^+$ diagonal con entradas
$(\boldsymbol{\Sigma}^+)_{jj} = 1/\sigma_j$ para $\sigma_j > 0$
y $0$ en caso contrario.
Cuando $\mathbf{A}$ es de rango completo, $\mathbf{A}^+ = (\mathbf{A}^\top\mathbf{A})^{-1}\mathbf{A}^\top$.
\end{proposicion}

\begin{proposicion}{Número de condición vía SVD}
\label{prop:cond_svd}
El número de condición espectral de $\mathbf{A}$ se calcula directamente
de la SVD como
\[
  \kappa_2(\mathbf{A}) \;=\; \frac{\sigma_{\max}}{\sigma_{\min}} \;=\; \frac{\sigma_1}{\sigma_r},
\]
donde $r$ es el rango numérico de $\mathbf{A}$ y $\sigma_{\min} = \sigma_r$
es el menor valor singular \emph{no nulo}. Esta es la forma computacionalmente
recomendada de calcular $\kappa(\mathbf{A})$, pues evita calcular
$\mathbf{A}^{-1}$ explícitamente.
\end{proposicion}

\begin{fisicaconexion}{SVD como análisis modal}
La SVD es el análogo algebraico de la descomposición modal en mecánica:
$\mathbf{A} = \sum_{k=1}^r \sigma_k\, \mathbf{u}_k \mathbf{v}_k^\top$.
Cada término de rango uno $\sigma_k\,\mathbf{u}_k\mathbf{v}_k^\top$ es un
``modo'' de la aplicación lineal, con ``amplitud'' $\sigma_k$. Truncar
la SVD a los primeros $\ell < r$ términos produce la mejor aproximación de
rango $\ell$ de $\mathbf{A}$ en norma de Frobenius (teorema de Eckart-Young).
En el contexto de la regularización espectral (Práctica 2), este truncamiento
es exactamente lo que hace el filtro de Tikhonov: atenúa los modos de
pequeño valor singular (ruidosos) mientras preserva los de valor singular
grande (señal).
\end{fisicaconexion}

%═══════════════════════════════════════════════════════════════════════════════
\section{Estimación empírica de la tasa de convergencia}
\label{ap:convergencia}
%═══════════════════════════════════════════════════════════════════════════════

En experimentos numéricos se trabaja frecuentemente con una secuencia de
soluciones calculadas con parámetro de discretización $h$ (paso de malla,
número de puntos, orden del método) y se desea estimar la tasa de
convergencia hacia la solución exacta.

\begin{definicion}{Tasa de convergencia algebraica}
\label{def:tasa_conv}
Se dice que un método converge con \textbf{tasa algebraica} $\alpha > 0$
si el error satisface
\[
  E(h) \;\sim\; C\,h^\alpha \quad \text{cuando } h \to 0,
\]
para alguna constante $C > 0$. Tomando logaritmos:
$\log E \approx \log C + \alpha \log h$.
La tasa se estima empíricamente a partir de dos experimentos con parámetros
$h_1$ y $h_2$ como
\[
  \hat{\alpha} \;=\; \frac{\log(E_1/E_2)}{\log(h_1/h_2)}.
\]
\end{definicion}

\begin{definicion}{Tasa de convergencia espectral (exponencial)}
\label{def:tasa_exp}
Cuando el error decrece más rápido que cualquier potencia de $h$,
se habla de convergencia \textbf{espectral} o exponencial:
\[
  E(p) \;\sim\; C\,e^{-\gamma p} \quad \text{cuando } p \to \infty,
\]
donde $p$ es el número de modos (grado o número de funciones de base).
En escala semilogarítmica, la curva $\log E$ vs.\@ $p$ es lineal con
pendiente $-\gamma$. Las bases de Fourier para funciones periódicas
analíticas exhiben esta convergencia, mientras que los monomios en
puntos equiespaciados presentan divergencia para muchas funciones
suaves (fenómeno de Runge).
\end{definicion}

\begin{ejemplo}{Estimación de $\hat{\alpha}$ en una tabla de resultados}
Suponga que al variar el número de puntos $n$ se obtienen los errores
$E_1 = 0.12$ con $n_1 = 10$ y $E_2 = 0.031$ con $n_2 = 20$ (equivalente
a $h_1 = 1/10$, $h_2 = 1/20$). Entonces
\[
  \hat{\alpha} = \frac{\log(0.12/0.031)}{\log(10/20)}
  = \frac{\log 3.87}{\log 0.5}
  = \frac{1.353}{-0.693} \approx 1.95,
\]
consistente con convergencia de segundo orden ($\alpha = 2$).
\end{ejemplo}

\begin{ejerciciocomp}{Verificación de la tasa de convergencia en mínimos cuadrados}
\label{ej:convergencia}
\dificultad{2}{1}
Considere $f(x) = \sin(2\pi x)$ en $[0,1]$ y ajústela con una base de
Fourier truncada a $p$ términos, con $n = 200$ puntos de entrenamiento.
Calcule el RMSE en un conjunto de prueba de $500$ puntos para
$p \in \{1, 3, 5, \ldots, 29\}$. Grafique $\log(\mathrm{RMSE})$ vs.\ $p$
y estime la tasa de convergencia espectral $\hat{\gamma}$. Discuta si la
convergencia es consistente con la regularidad de $f$.
\end{ejerciciocomp}

%═══════════════════════════════════════════════════════════════════════════════
\section{Guía de referencia rápida}
\label{ap:guia_rapida}
%═══════════════════════════════════════════════════════════════════════════════

La \Cref{tab:referencia_rapida} consolida las fórmulas de uso más frecuente
en las prácticas, organizadas por la tarea computacional que realizan.

\begin{table}[htbp]
\centering
\caption{Referencia rápida de fórmulas computacionales. SVD thin:
  $\mathbf{A} = \hat{\mathbf{U}}\hat{\boldsymbol{\Sigma}}\mathbf{V}^\top$
  con $\hat{\mathbf{U}} \in \mathbb{R}^{n\times p}$,
  $\hat{\boldsymbol{\Sigma}} \in \mathbb{R}^{p\times p}$,
  $\mathbf{V} \in \mathbb{R}^{p\times p}$.}
\label{tab:referencia_rapida}
\small
\begin{tabular}{lll}
  \toprule
  \textbf{Tarea} & \textbf{Fórmula} & \textbf{Notas} \\
  \midrule
  Solución por eq.\ normales
    & $\boldsymbol{\theta}^* = (\mathbf{A}^\top\mathbf{A})^{-1}\mathbf{A}^\top\mathbf{y}$
    & Rango completo; inestable si $\kappa(\mathbf{A})^2$ grande \\[3pt]
  Solución por SVD
    & $\boldsymbol{\theta}^* = \mathbf{V}\hat{\boldsymbol{\Sigma}}^{-1}\hat{\mathbf{U}}^\top\mathbf{y}$
    & Estable; preferida computacionalmente \\[3pt]
  Número de condición
    & $\kappa_2(\mathbf{A}) = \sigma_1 / \sigma_p$
    & Vía SVD thin; $\sigma_1 \geq \cdots \geq \sigma_p > 0$ \\[3pt]
  RMSE
    & $\mathrm{RMSE} = \|\mathbf{r}\|_2 / \sqrt{n}$
    & $\mathbf{r} = \hat{\mathbf{y}} - \mathbf{y}$ \\[3pt]
  Error máximo
    & $E_\infty = \|\mathbf{r}\|_\infty$
    & Sensible a outliers \\[3pt]
  MAE
    & $\mathrm{MAE} = \|\mathbf{r}\|_1 / n$
    & Robusto frente a outliers \\[3pt]
  Tasa de convergencia
    & $\hat{\alpha} = \log(E_1/E_2)/\log(h_1/h_2)$
    & Estimador de dos puntos \\[3pt]
  Cond.\ ecuaciones normales
    & $\kappa(\mathbf{A}^\top\mathbf{A}) = \kappa(\mathbf{A})^2$
    & Siempre usar SVD en su lugar \\[3pt]
  Pseudoinversa
    & $\mathbf{A}^+ = \mathbf{V}\hat{\boldsymbol{\Sigma}}^{-1}\hat{\mathbf{U}}^\top$
    & Mínima norma; generaliza inv.\ para rect. \\
  \bottomrule
\end{tabular}
\end{table}

\FloatBarrier

\begin{notasbib}
El tratamiento estándar de referencia para normas matriciales y número de
condición es \textcite{golub2013matrix}, que también cubre la SVD con
detalle algorítmico completo. Para análisis de estabilidad y efectos de
punto flotante en álgebra lineal numérica, la fuente definitiva es
\textcite{higham2002accuracy}; en particular, los Capítulos 2 y 7 son
directamente relevantes para las prácticas de este laboratorio. La
interpretación del número de condición como razón de amplificación de
perturbaciones sigue la presentación de \textcite{trefethen1997numerical}
(Lectures 12 y 14), texto altamente recomendado por su enfoque geométrico.
Para espacios $L^2$ y funciones de base, el marco funcional riguroso se
encuentra en \textcite{reed1980functional}. La estimación empírica de tasas
de convergencia sigue el capítulo introductorio de \textcite{leveque2007finite}.
\end{notasbib}

\printbibliography[heading=subbibliography,title={Referencias del Apéndice A}]